\documentclass[../../../include/open-logic-section]{subfiles}

\begin{document}
\olfileid{sth}{ordinals}{zfm}
\olsection{$\ZFminus$: एक महत्त्वाचा टप्पा}

प्रतिस्थापनाला समर्थन कसे द्यावे, किंवा देता येईलच का, हा प्रश्न सरळ नाही. म्हणून तो \olref[replacement][]{chap} साठी राखून ठेवू. पण प्रतिस्थापन जोडल्यामुळे आपण आणखी एक महत्त्वाचा टप्पा गाठला आहे. आता $\ZFminus$ या उपपत्तीसाठी लागणारी सर्व स्वयंसिद्धके आपल्याकडे आहेत. ती अशी:

\begin{defn}
$\ZFminus$ उपपत्तीत पुढील स्वयंसिद्धके आहेत: विस्तारात्मकता, संयोग, जोड्या, घातसंच, अनंतता, आणि विभक्तीकरण व प्रतिस्थापन या रूपबंधांतील सर्व स्वयंसिद्धके. दुसऱ्या शब्दांत, $\Zminus$ मध्ये प्रतिस्थापन जोडले की $\ZFminus$ मिळते.
\end{defn}
\noindent
यातील नाव \emph{झर्मेलो--फ्रेंकेल} संच उपपत्ती दर्शवते (त्यातून \emph{वगळलेली} गोष्ट आपण पुढे पाहू). प्रतिस्थापनाची मांडणी \citeyear{Fraenkel1922} मध्ये फ्रेंकेल यांनी केली असे मानले जाते; म्हणून त्यांच्या नावाला येथे स्थान आहे. मात्र त्याचे पहिले काटेकोर स्वरूप \citet{Skolem1922} यांनी दिले.

\end{document}
